\documentclass{beamer}
\input{tbook-deck}
\title{When fractions turn into integers}
\subtitle{Hanna's A181074 and moving exponents}
\author{Paper by Haobo Ma and Wenlin Zhang}
\institute{ChronoAI Pte Ltd, Singapore\\National University of Singapore (Wenlin Zhang)}
\date{Source: version 3, DOI 10.5281/zenodo.24408311}
\begin{document}
% Manuscript: Section 1 and A181074, Comments (October 2, 2010)
\begin{frame}
  \titlepage
  \note{Credit Paul D. Hanna for the question and Haobo Ma and Wenlin Zhang for the paper. This is a reading deck; the source author line does not name a live presenter.}
\end{frame}
\begin{frame}{Outline}
  \tableofcontents
  \note{Start with a coefficient anyone can compute. Then change the representation and show how the arithmetic survives changes to the exponent.}
\end{frame}
\section{The puzzle}
% Manuscript: Section 1, equation (1); Corollary 6, A181074 clause
\begin{frame}{Two sources of denominators}
  \begin{bookresult}{a181074-exact}
  Every coefficient of
  \[
    A(x)=\exp\!\left(\sum_{n\ge1,k\ge0}
        \binom{n+k-1}{k}^{k+1}\frac{x^{n+k}}n\right)
  \]
  is an integer.
  \end{bookresult}
  \note{The factors one over n and one over factorial appear in different steps. Degree m only involves the pairs with n plus k at most m; nothing depends on analytic convergence.}
\end{frame}
% Manuscript: Equation (1), cubic truncation; Section 7 initial terms
\begin{frame}{Why the cubic coefficient is five}
  \[
  \begin{array}{c|ccc}
    (n,k) &(1,2)&(2,1)&(3,0)\\\hline
    \text{contribution to }[x^3]L&1&2&1/3
  \end{array}
  \]
  \[L=x+\tfrac32 x^2+\tfrac{10}{3}x^3+\cdots,\qquad
    [x^3]\exp L=\tfrac{10}{3}+\tfrac32+\tfrac16=5.\]
  \note{The cross term in L squared divided by two is three halves. Checking more degrees still only checks finitely many. Ask what a denominator-free representation would look like.}
\end{frame}
\section{A denominator-free formula}
% Manuscript: Section 3, Lemma 4 and the Euler product
\begin{frame}{Factor before you expand}
  \[(1-x)^{-1}(1-x^2)^{-1}(1-x^3)^{-3}
      =1+x+2x^2+5x^3+\cdots.\]
  \[g_m=m[x^m]L,\qquad
       h_m=\frac1m\sum_{d\mid m}\mu(d)g_{m/d}.\]
  If every $h_m$ is integral, $\prod_m(1-x^m)^{-h_m}$ is integral.
  \note{The first ghost coefficients are one, three and ten. The inversion formula gives h one, h two and h three equal to one, one and three. Negative h is harmless because its factor is a polynomial.}
\end{frame}
% Manuscript: Section 3, proof of Lemma 4
\begin{frame}{Pair the prime-scaled diagonals}
  \[g_m=\sum_{n=1}^m \frac{m}{n}F(n,m-n).\]
  \begin{itemize}
    \item $n\mid kF(n,k)$ makes each summand integral.
    \item In $g_{pm}$, discard the terms with $p\nmid n$ modulo $p^{v_p(m)+1}$.
    \item Pair $n=pj$ with $j$ in $g_m$ and use the array congruence.
  \end{itemize}
  \note{After pairing, the array difference contributes v p of j plus one, and m over j changes the valuation to v p of m plus one. In the Mobius sum, pair divisors d and p times d to prove divisibility by every prime power in m.}
\end{frame}
\section{Moving exponents and their limits}
% Manuscript: Theorem 1 and Lemmas 2--3
\begin{frame}{Two ways to preserve a residue}
  \begin{bookresult}{general-integrality}
  For $U\in1+t\mathbb Z[[t]]$, the exponential built from
  $([t^k]U^n)^E$ is integral for $E=rk+s$ ($r,s\ge1$),
  $E=rn+sk+q$ ($r\ge1$, $s,q\ge0$), or $E=n+k-1$.
  \end{bookresult}
  \begin{itemize}
    \item Prime-to-$p$ base: the exponent difference is a multiple of the period.
    \item Base divisible by $p$: both powers vanish modulo $p^{v_p(n)+1}$.
  \end{itemize}
  \note{Lift the coefficient first. For k zero the base is one, including the unique possible zero exponent. For a nonzero nonunit base, the power contributes the missing copies of p. The congruence and divisibility then feed the diagonal argument.}
\end{frame}
% Manuscript: Corollaries 5, 9 and 11; Section 7
\begin{frame}{What changes when we change the array?}
  \begin{itemize}
    \item For a degree bound $d$, use distance $dn-k$ from the opposite edge.
    \item Keep $x^ny^k$ rather than $x^{n+k}$ to retain a triangle.
    \item Lower $k+1$ to $k$ in A181074: $[x^6]A=1481/2$ for the altered series.
  \end{itemize}
  \note{The reflected polynomial statement allows negative coefficients and degrees below the bound. The bivariate criterion needs finite support per row. The displayed fraction is a counterexample for the altered series, not a failure of Hanna's original conjecture.}
\end{frame}
% Manuscript: Section 1, Sections 5--7, references
\begin{frame}{What the argument adds}
  \begin{itemize}
    \item Moree and Gossow: divisibility and congruences before the moving power.
    \item Du--Huang--Li and Byszewski--Graff--Ward: Dold and Euler-product criteria.
    \item Ma and Zhang: preserve the congruences through changing exponents.
  \end{itemize}
  \note{The companion root limits use a different argument, comparing largest terms via Stirling estimates. Finite numerical checks are examples, not the proof. The all-degree A181074 claim and the general families have different formal evidence.}
\end{frame}
\appendix
% Manuscript: Verification ledger (generated)
\FormalVerification
\end{document}
